\chapter{Physical Security Breaches}
\label{ch:2.2}
\chaptertopic{Topic 2.2}

\chapterquote{The only truly secure system is one that is powered off, cast in a block of concrete and sealed in a lead-lined room with armed guards --- and even then I have my doubts.}{Eugene H. Spafford, quoted in \emph{Scientific American} (March 1989)}

\section{Motivation}

Chapter~\ref{ch:2.1} gave you the vocabulary to reason about any security goal in the abstract.
This chapter grounds that vocabulary in the oldest attack surface there is: the physical building
and the physical device. Firewalls and encryption dominate the popular image of cybersecurity. But a remarkable share of
real breaches begin with an attacker simply walking through a door, or with a laptop that leaves
an office and comes back subtly different. Locks, badge readers, HVAC systems, and disk encryption
all defend the same underlying asset that Chapter~\ref{ch:2.1} described abstractly. Here, though,
the threats are a held-open door, a burst pipe, and a stolen hard drive, not a phishing email.

\begin{apalignment}{Unit 2: Securing Spaces}{Topic 2.2 -- Physical Vulnerabilities and Attacks}
\begin{apobjectives}
\apobjective{2.2.A}{Analyze Risk}{Identify the physical attacks an organization is realistically likely to face, from following someone through a controlled door to defeating a lock, duplicating a credential, or removing equipment outright.}
\apobjective{2.2.B}{Analyze Risk}{Explain the causal path by which a threat exploits a physical weakness and produces a specific outcome for an asset -- its loss, its damage, the disruption of the service it supports, or its destruction.}
\apobjective{2.2.C}{Analyze Risk}{Assess the risk arising from an identified physical weakness and record it in a form that supports a decision: what was observed, what it would allow, how likely and how costly, and who owns the outcome.}
\end{apobjectives}
\end{apalignment}

\section{Unauthorized Physical Access}

\begin{definition}{Tailgating}
Following an authorized person through a secured door without using one's own credentials, while
the authorized person is unaware, or does not consent to, the attacker's entry.
\end{definition}

Tailgating exploits an ordinary social norm: holding a door for the person behind you. An attacker
carrying a box or coffee, hands visibly full, approaches a card-reader door timed just behind a
legitimate employee. The employee swipes in and holds the door out of courtesy, and the attacker
walks through having scanned no badge at all. The access log shows only the legitimate employee's
entry -- the attacker leaves no direct trace in the system meant to record exactly this event.
This is the same social pressure studied in Chapter~\ref{ch:1.1}: urgency, authority cues, and
plain politeness make people reluctant to demand a badge from someone who looks like they belong.

\begin{definition}{Piggybacking}
The consensual counterpart to tailgating: an authorized employee knowingly lets an unverified
person through a secured entry, usually as a favor to a familiar contractor or visitor.
\end{definition}

In an access log, piggybacking is indistinguishable from tailgating -- one badge swipe, two
people through the door. The difference is intent, but the security exposure is identical: an
unverified person now has physical access. Recognition is not authorization; a familiar face may
be a former employee whose access was supposed to be revoked.

When tailgating fails or is impractical, an attacker may target the lock itself. \textbf{Lock
picking} manipulates a pin-tumbler lock's pins with specialized tools, requiring only hours of
practice against simple locks and needing no cryptographic knowledge at all. A \textbf{bump key}
is cut to a generic profile and, struck with a mallet while turned, momentarily bounces every pin
past the shear line at once. \textbf{Shimming} slides a thin metal shim into a padlock's shackle to
retract its locking mechanism directly. And \textbf{RFID badge cloning} captures the radio signal
of an access badge with a concealed reader and replays it from a blank credential. This is trivial
against older 125~kHz proximity cards, which carry no cryptographic protection at all, with
hardware costing under fifty dollars.

\begin{example}[Physical access as a force multiplier]
Once an attacker reaches a device physically, most software-layer controls stop mattering. A
locked laptop can be booted from an external USB drive to mount its disk directly or reset its
local administrator password. Either route entirely bypasses the operating-system login screen it
was never asked to authenticate against. An Ethernet port behind a bypassed lock can host a rogue
network tap. Full-disk encryption is one of the very few controls that survives physical
possession. Even then, it only holds if the device was powered off when it was taken, so its
decryption key was not sitting in memory waiting to be read.
\end{example}

Access logs can still surface these attacks algorithmically, even without a human watching every
door. A well-instrumented system flags a physical impossibility: one badge swiping door D1 at
08:04:11 and door D2 (some distance away) at 08:04:13 did not happen by one person walking
normally. It also flags a door logging a single swipe while a separate camera sensor counts two
people crossing -- that mismatch is tailgating by definition, whether or not anyone consented.

\begin{lstlisting}[language=Python, caption={Flagging tailgating and impossible-speed badge events}]
def flag_access_anomalies(events, min_seconds_between_doors=10):
    """events: list of dicts like {"badge_id": str, "door_id": str,
       "timestamp": int, "crossings_detected": int}, sorted by timestamp
       is not assumed. Returns sorted event indices flagged as anomalous:
       either a door recorded more physical crossings than badge swipes,
       or the same badge appears at two different doors too quickly."""
    flagged = set()
    for i, e in enumerate(events):
        if e["crossings_detected"] > 1:
            flagged.add(i)

    by_badge = {}
    for i, e in enumerate(events):
        by_badge.setdefault(e["badge_id"], []).append((e["timestamp"], i, e["door_id"]))
    for badge_events in by_badge.values():
        badge_events.sort()
        for (t1, i1, d1), (t2, i2, d2) in zip(badge_events, badge_events[1:]):
            if d1 != d2 and t2 - t1 < min_seconds_between_doors:
                flagged.add(i1)
                flagged.add(i2)
    return sorted(flagged)
\end{lstlisting}

The corresponding defenses target the exploit, not the symptom. A \textbf{mantrap} -- two
interlocked doors where the second cannot open until the first closes -- removes the physical
opportunity for a second person to follow through at all. Upgrading legacy proximity cards to
encrypted, mutually authenticated 13.56~MHz cards (MIFARE DESFire) removes the cloning
vulnerability. And awareness training built around a simple rule -- one badge, one person -- turns
the social norm that tailgating exploits into a norm that defeats it instead.

\begin{misconception}
\textbf{``A locked door means the room is secure.''} A lock is a delay mechanism, not an absolute
barrier. Lock-bypass techniques, credential cloning, and tailgating all defeat it without leaving
the kind of obvious evidence a broken window would.
\end{misconception}

\section{Environmental Vulnerabilities}

Not every threat to a server room is a person trying to get in. A server room can be destroyed
just as thoroughly by conditions that are never ``unauthorized'' in the access-control sense at
all. Because these are not attacks in the social-engineering sense, they are easy to
under-prioritize -- even though the outcome is the same as ransomware: systems go offline, and
sometimes data is destroyed outright.

\textbf{Power.} A total outage stops everything immediately absent backup power, and a merely
``dirty'' supply -- voltage spikes, sags, brownouts -- can corrupt in-flight writes or degrade
power supplies over time. An \textbf{uninterruptible power supply (UPS)} bridges the gap between
an outage and generator start, or rides out a short dip entirely. Data centers commonly
design for \textbf{N+1 redundancy}: enough generator and UPS capacity that any single unit can
fail without causing an outage. A generator that has not been load-tested in a year is not a
control at all, just an assumption.

\textbf{HVAC.} Cooling is not a comfort feature in a server room; it is life support for hardware
that throttles or shuts down once internal temperatures cross a safety threshold. A broken chiller
or a clogged filter can push rack temperatures past safe limits within minutes in a dense room.
Humidity matters in both directions -- too little raises static-discharge risk to circuit boards,
too much promotes condensation and corrosion. That is why facilities target roughly
18--27\textdegree{}C and 40--60\% relative humidity per ASHRAE guidelines, alarming on either
threshold.

\textbf{Fire.} Fire is the most destructive environmental threat, and a server room complicates
the standard response. Water-based sprinklers, the default suppressant almost everywhere else,
will destroy the exact equipment they are meant to protect. Data centers instead use
\textbf{clean-agent fire suppression} (FM-200, inert-gas systems) that extinguishes fire by
displacing oxygen or interrupting the chemical reaction without damaging electronics. This is
paired with early smoke and heat detection tuned to trigger well before a sprinkler-grade
threshold would.

\textbf{Water and flooding.} Burst pipes, groundwater intrusion in basement facilities, and roof
leaks all threaten electronics through short-circuiting and corrosion. Burst pipes are especially
dangerous when plumbing is routed above a server room -- a design flaw whenever a data center
shares a building with restrooms or kitchens upstairs. \textbf{Raised flooring} elevates equipment
above the slab so a shallow flood does not reach it directly. Floor-level water sensors trigger an
alarm at the first sign of moisture, rather than after damage has already occurred.

\begin{lstlisting}[language=Python, caption={Triaging environmental sensor readings by severity}]
def triage_reading(temp_c, humidity_pct, smoke_detected, water_detected):
    """Returns "critical", "warning", or "normal" for one sensor snapshot."""
    if smoke_detected or water_detected:
        return "critical"
    temp_bad = not (18 <= temp_c <= 27)
    humidity_bad = not (40 <= humidity_pct <= 60)
    if temp_bad and humidity_bad:
        return "critical"
    if temp_bad or humidity_bad:
        return "warning"
    return "normal"
\end{lstlisting}

All four categories -- power, HVAC, fire, water -- are first and foremost availability threats. An
outage from any one of them takes systems offline the same way a DDoS attack does, except the
``attacker'' is a physical process rather than a person. A mature risk analysis, using
Chapter~\ref{ch:2.1}'s vocabulary, scores these risks with the same rigor as a targeted attack,
because the impact to the organization can be identical or worse.

\begin{misconception}
\textbf{``Sprinklers are fine for a server room like anywhere else in the building.''}
Water-based sprinklers actively destroy electronics; a data center's fire suppression is
specifically chosen so that a false alarm, or even a real fire, does not itself become a
total-loss event.
\end{misconception}

\section{Device Theft and Tampering}

Unauthorized access gets an attacker into a building. Once a device is reachable -- or once it
leaves the building in an employee's bag -- the attacker's goal shifts to the device itself.

\begin{definition}{Device tampering}
A covert, typically brief physical alteration of a device -- hardware, firmware, or a peripheral
-- intended to compromise it without the owner's knowledge, as distinct from device theft, which
is the device's outright physical removal.
\end{definition}

The naive framing of theft is ``we lost a \$1{,}500 laptop''; the real cost is almost always the
data, not the hardware. A stolen \emph{unencrypted} laptop hands an attacker every file, cached
credential, and saved browser session on the disk, often without guessing a single password.
Booting from external media lets an attacker mount the disk or reset the local administrator
account without ever going through the intended login screen. \textbf{Full-disk encryption} is the
single most effective control against this specific threat. If the drive is encrypted and the
device was off -- so its key was never resident in memory -- when it was taken, the stolen disk is
unreadable without the key.

\begin{example}[The evil-maid attack]
An attacker with brief, unsupervised physical access to a device -- the textbook case is a hotel
room cleaner -- boots a malicious USB image that replaces the boot loader with a fake password
prompt visually identical to the real one. The attacker then leaves. The next time the legitimate
owner enters their disk-encryption password, the fake loader silently captures it and reboots into the genuine
system as though nothing happened. The owner never notices anything wrong; the attacker now holds
the encryption key and can return later, or exfiltrate it over a network channel, to unlock the
data at leisure. This defeats full-disk encryption not by breaking any cryptography, but by
exploiting the brief unsupervised window \emph{before} that session's encryption was ever engaged.
\end{example}

\begin{figure}[h]
\centering
\begin{tikzpicture}[
  evilstep/.style={draw=kickstartermaroon, thick, rounded corners, minimum width=3.6cm, minimum height=1cm, align=center, font=\scriptsize},
  yshift/.style={},
]
\node[evilstep, fill=kickstartercream] (t1) at (0,0) {1. Attacker gets brief\\unsupervised access};
\node[evilstep, fill=kickstartermaroon!20] (t2) at (4.1,0) {2. Boots malicious USB,\\installs fake password\\prompt, leaves};
\node[evilstep, fill=kickstartercream] (t3) at (8.2,0) {3. Owner returns, types\\real password into\\the fake prompt};
\node[evilstep, fill=kickstartermaroon!20] (t4) at (4.1,-2) {4. Fake loader captures\\the key, reboots into\\the genuine system};
\node[evilstep, fill=kickstartercream] (t5) at (8.2,-2) {5. Owner sees nothing\\wrong; attacker holds\\the encryption key};
\draw[-{Latex[length=2mm]}] (t1) -- (t2);
\draw[-{Latex[length=2mm]}] (t2) -- (t3);
\draw[-{Latex[length=2mm]}] (t3) -- (t4);
\draw[-{Latex[length=2mm]}] (t4) -- (t5);
\node[font=\scriptsize\sffamily, gray] at (4.1,1) {the cryptography is never broken -- only the moment before it engages};
\end{tikzpicture}
\caption{The evil-maid attack succeeds entirely in the window before disk encryption is engaged
for the session; the encryption itself is never attacked.}
\label{fig:2.2-evil-maid}
\end{figure}

Tampering is not limited to software. A rogue USB implant placed inline between a keyboard and a
computer silently logs every keystroke, including passwords, and can exfiltrate over its own
hidden radio. A ``free'' USB charging cable or promotional drive can present itself to the
operating system as an innocuous device while secretly acting as a keyboard that types malicious
commands the moment it is plugged in. Because these implants are physically small and often
hidden in plain sight, tampering can go undetected indefinitely unless staff are specifically
trained to notice inconsistencies.

Detection here relies on physical evidence rather than a login prompt, since tampering is
engineered to leave no digital trace. \textbf{Tamper-evident seals} -- stickers that visibly tear
or reveal a ``VOID'' pattern if removed -- over a laptop's screws or a server's chassis panel let
staff visually confirm a device has not been opened. Scheduled comparison of serial numbers and
asset tags against inventory catches a swapped device. And staff trained to notice an unfamiliar
dongle in a keyboard port, or a boot that took slightly longer than usual, generate the kind of
soft signal that should trigger a review rather than be dismissed.

\begin{lstlisting}[language=Python, caption={Triaging a device incident report by organizational risk}]
def triage_device_incident(encrypted, seal_intact, hours_since_last_check):
    """Returns "LOW", "MEDIUM", or "HIGH" risk for a reported theft/
    tampering incident."""
    if not encrypted and hours_since_last_check > 24:
        return "HIGH"
    if not seal_intact:
        return "HIGH" if not encrypted else "MEDIUM"
    if not encrypted:
        return "MEDIUM"
    return "LOW"
\end{lstlisting}

Organizations reduce the impact of theft with asset tags and inventory tracking (so a missing
device is noticed in hours, not months), \textbf{remote wipe} capability through mobile-device
management tools, and cable locks that raise the bar against opportunistic theft from a shared
workspace. None of these stop a prepared, determined attacker, and none of them address tampering
at all, since tampering never requires removing the device from the room.

\begin{misconception}
\textbf{``If a laptop is password-protected, a thief can't get the data.''} A login password
protects the normal operating-system boot path only. An attacker with physical possession can
usually boot from external media and read the disk directly unless the disk's contents are
themselves encrypted.
\end{misconception}

\begin{example}[Writing a physical risk register entry that actually supports a decision]
``Building security is inadequate'' cannot be acted on -- it names no observation, no mechanism,
and no owner. A useful entry names all four things the treatment decision in Chapter~\ref{ch:2.1}
needs. It states what was observed: the loading-dock door was found propped open during three of
five unannounced site visits. It states what that would allow: an unbadged person can reach the
receiving corridor and, from there, an unlocked electrical room. It gives an estimate of
likelihood and impact: moderate likelihood given the recurring pattern, high impact given the
electrical room's contents. And it names an owner accountable for the finding until it is closed:
the facilities manager, with a documented remediation date. A risk register entry that skips any
one of these four elements -- especially the owner -- tends to stay open indefinitely, because no
specific person is answerable for it.
\end{example}

\section*{Review Questions}
\begin{reviewquestions}
\item An access log shows door D1 recording a single badge swipe while a separate camera sensor
counts two people entering within the same second. Explain whether this is tailgating or
piggybacking, and explain why the distinction does not change the security response.
\item Explain why a locked, powered-on laptop with full-disk encryption enabled is still
vulnerable to an evil-maid attack, and identify the specific moment in the attack where the
encryption is defeated.
\item A facility's sensors report rising temperature and rising humidity simultaneously. Explain
why this combination is treated as more severe than either reading alone, and name the underlying
failure it most likely indicates.
\item Explain why data centers use clean-agent fire suppression instead of water sprinklers, and
identify which CIA pillar (from Chapter~\ref{ch:2.1}) this control is primarily protecting.
\item A device incident report shows an unencrypted laptop with an intact tamper seal, last
inventoried 30 hours ago. Using the triage logic from this chapter, determine its risk level and
explain which single factor is doing the most work in that classification.
\item In the evil-maid attack, the disk encryption itself is never broken. Explain precisely what
the attacker does capture, and why a tamper-evident seal over the laptop's chassis would not have
prevented this specific attack even though it would have revealed physical tampering with the
hardware.
\item A risk register entry reads only ``building security is inadequate.'' Rewrite it in your
own words to include the four elements this chapter requires, using a physical weakness of your
choosing. Then explain why omitting the owner in particular tends to leave a finding open
indefinitely.
\end{reviewquestions}

\begin{answerkey}
\item This is tailgating or piggybacking depending on intent, which the log alone cannot
distinguish -- both produce exactly one badge swipe against two people crossing. The distinction
does not change the security response because the exposure is identical in either case: an
unverified person now has physical access to the space. The appropriate fix -- a mantrap,
awareness training, or an anti-passback rule flagging the crossing-count mismatch -- targets the
physical opportunity for a second person to pass through, not the second person's motive.
\answerlabel{2.2.A}
\item Full-disk encryption only protects data once the encryption key is not resident anywhere an
attacker can read it. An evil-maid attack defeats it not by breaking the cryptography but by
exploiting the brief unsupervised window \emph{before} the owner next enters their password. The
attacker replaces the boot loader with a fake, visually identical prompt, and the owner's next
login hands the real key to the attacker's captured prompt rather than to the genuine decryption
process. The encryption is defeated at the moment the owner types their password into the
attacker's fake loader, not at any later point.
\answerlabel{2.2.B}
\item Rising temperature and rising humidity together, rather than either alone, typically indicate
an HVAC failure rather than a localized problem such as one failing fan or a single humidifier
malfunction. This is because a working HVAC system actively regulates both variables
simultaneously, and only a systemic failure lets both drift out of range together. The combination
is treated as more severe because it signals a facility-wide root cause requiring immediate
escalation to facilities management, rather than a narrow fix like replacing one component.
\answerlabel{2.2.B}
\item Water-based sprinklers, the default fire suppressant elsewhere, would themselves destroy the
electronics they were meant to save, converting a fire event into a guaranteed total-loss event for
every device in the room. Clean-agent systems (FM-200, inert gas) suppress the fire by displacing
oxygen or interrupting the combustion reaction without depositing any conductive or corrosive
residue on hardware. This control primarily protects \textbf{availability}: its entire purpose is
keeping systems reachable and operational, not confidential or unaltered.
\answerlabel{2.2.C}
\item Applying the triage function: \texttt{encrypted=False}, \texttt{seal\_intact=True},
\texttt{hours\_since\_last\_check=30}. Because the device is unencrypted and more than 24 hours have
elapsed since the last check, the first branch fires and the incident is classified \textbf{HIGH}
risk. The factor doing the most work is the missing encryption combined with the extended
uninventoried window: an encrypted device in the same scenario would only reach MEDIUM, because an
attacker without the key gains nothing from physical possession alone. An unencrypted device left
unaccounted for over a day, by contrast, gives an attacker ample time to extract everything on the
disk.
\answerlabel{2.2.C}
\item The attacker captures the decryption password itself, typed by the owner into a fake boot
loader that is visually indistinguishable from the real one -- not any weakness in the encryption
algorithm. A tamper-evident seal over the chassis would not prevent this because the evil-maid
attack does not require opening the case at all. The attacker boots the machine from an external
USB drive through the normal boot process, replaces the boot loader in software, and closes the
lid, leaving every physical seal untouched. The seal detects hardware-level intrusion; this attack
is entirely a software substitution that happens between two ordinary boots.
\answerlabel{2.2.B}
\item A usable entry names all four elements. What was observed: for example, the server-room
door badge reader has logged three tailgating-pattern crossing-count mismatches in the past
month. What it would allow: an unbadged individual reaching racks containing production database
servers. Likelihood and impact: moderate likelihood given the recurring pattern, high impact given
the asset. And a named owner accountable until the finding closes: the facilities or
physical-security manager, with a documented remediation date. Omitting the owner tends to leave a
finding open indefinitely, because a risk with no specific person answerable for it has no one
whose job performance is affected by it staying unresolved. It becomes everyone's problem in
theory and no one's in practice.
\answerlabel{2.2.C}
\end{answerkey}
